% Qualitative effect of QLoRA fine-tuning — label consistency evidence
% Card: test_germany_0015.jpg
\begin{table}[h]\centering
\caption{Qualitative Effect of QLoRA Fine-tuning on VLM Output Format}
\label{tab:finetune_qual}
\begin{tabular}{lll}
\toprule
Field & Base Model Output & Fine-tuned Output \\
\midrule
Name      & \texttt{Name | Leonie D\"{o}rr}         & \texttt{Name | Leonie D\"{o}rr} \\
DOB       & \texttt{Date of Birth | 01.08.1991}    & \texttt{Date of Birth | 01.08.1991} \\
ID Number & \texttt{\textit{[empty]} | M898585166} & \texttt{ID No. | M898585166} \\
Signature & \texttt{Signature | [none]}             & \texttt{Signature | [none]} \\
\bottomrule
\end{tabular}
\vspace{2pt}\\
\footnotesize{The base model detects PII values correctly but omits field labels on
structured fields (row~3). Fine-tuning teaches consistent \texttt{LABEL\,|\,VALUE}
output format required by the downstream pipeline: an empty label is normalised to
\texttt{other} by \texttt{normalize\_cat()} and scored as a false positive regardless
of whether the value is correct. Quantitative F1 appears identical between the two
models (Table~\ref{tab:finetune}) because the evaluation framework matches on values,
not on the label\,$\rightarrow$\,category mapping the live pipeline depends on.}
\end{table}
